\documentclass[12pt,a4paper]{article}
\usepackage[utf8]{inputenc}
\usepackage[spanish]{babel}
\usepackage{graphicx}
\usepackage{hyperref}
\usepackage{geometry}
\usepackage{fancyhdr}
\usepackage{listings}
\usepackage{xcolor}
\usepackage{titlesec}

% Configuración de página
\geometry{margin=2.5cm}

% Configuración de enlaces
\hypersetup{
    colorlinks=true,
    linkcolor=blue,
    filecolor=magenta,      
    urlcolor=cyan,
    pdftitle={Primer Avance - Todo App},
}

% Configuración de código
\lstset{
    basicstyle=\ttfamily\small,
    breaklines=true,
    frame=single,
    backgroundcolor=\color{gray!10},
    commentstyle=\color{green!50!black},
    keywordstyle=\color{blue},
    stringstyle=\color{red}
}

% Encabezado y pie de página
\pagestyle{fancy}
\fancyhf{}
\rhead{Primer Avance - Todo App}
\lhead{DevOps}
\rfoot{Página \thepage}

\begin{document}

% ===============================
% PORTADA
% ===============================
\begin{titlepage}
    \centering
    \vspace{2cm}
    
    {\Huge\bfseries Primer Avance del Proyecto\par}
    \vspace{1cm}
    {\LARGE Todo App\par}
    \vspace{0.5cm}
    {\large Aplicación Full Stack de Gestión de Tareas\par}
    
    \vspace{3cm}
    
    {\Large\textbf{Integrantes:}\par}
    \vspace{0.5cm}
    {\large
        JAVIER ALBEIRO AGAMEZ SILGADO\\[0.3cm]
        OSCAR DUVAN ALVAREZ CARO\\[0.3cm]
        NEIDY MICHELLE CRUZ ROSALEZ\\[0.3cm]
        MICHELLE MOREA NINCO\\[0.3cm]
        LEANDRO NICOLAS VARGAS CUBIDES\\[0.3cm]
    }
    
    \vfill
    
    {\large\textbf{Fecha de entrega:} 04 de Noviembre de 2025\par}
    \vspace{0.5cm}
    {\large\textbf{Repositorio GitHub:}\par}
    \vspace{0.3cm}
    {\small\url{https://github.com/LeandroNV/todo-app-fullstack}\par}
    
    \vspace{1cm}
\end{titlepage}

% ===============================
% ÍNDICE
% ===============================
\tableofcontents
\newpage

% ===============================
% DESCRIPCIÓN TÉCNICA
% ===============================
\section{Descripción Técnica}

\subsection{¿Qué se hizo en este avance?}

En este primer avance se desarrolló una aplicación completa de gestión de tareas (To-Do App) siguiendo la arquitectura de microservicios. El proyecto incluye:

\begin{itemize}
    \item \textbf{Frontend completo:} Desarrollado con Vue 3, TypeScript y Tailwind CSS, implementando 4 componentes principales (TodoForm, TodoItem, TodoFilters, TodoStats).
    
    \item \textbf{Backend robusto:} API REST completa con Node.js, Express y TypeScript, incluyendo 10 endpoints para operaciones CRUD y estadísticas.
    
    \item \textbf{Base de datos:} MongoDB con Mongoose como ODM, implementando validaciones y índices para optimizar el rendimiento.
    
    \item \textbf{Dockerización completa:} Se crearon Dockerfiles para cada servicio y un docker-compose.yml para orquestar todos los contenedores.
    
    \item \textbf{Documentación exhaustiva:} 6 archivos de documentación (README, SETUP, QUICKSTART, etc.) y scripts de utilidad.
\end{itemize}

\subsection{¿Qué servicios contiene el sistema?}

El sistema está compuesto por tres servicios principales:

\begin{enumerate}
    \item \textbf{Frontend (Vue 3)}
    \begin{itemize}
        \item Puerto: 80
        \item Tecnologías: Vue 3, TypeScript, Tailwind CSS, Vite
        \item Servidor: Nginx (en producción)
        \item Características: SPA responsive con estado reactivo
    \end{itemize}
    
    \item \textbf{Backend (Node.js + Express)}
    \begin{itemize}
        \item Puerto: 3000
        \item Tecnologías: Node.js 20, Express, TypeScript, Mongoose
        \item API REST con 10 endpoints
        \item Validaciones y manejo de errores centralizado
    \end{itemize}
    
    \item \textbf{Base de Datos (MongoDB)}
    \begin{itemize}
        \item Puerto: 27017
        \item Versión: MongoDB 7 (Alpine)
        \item Persistencia: Volúmenes Docker
        \item Healthcheck configurado
    \end{itemize}
\end{enumerate}

\subsection{¿Cómo se conectan?}

La arquitectura de conexión sigue el siguiente flujo:

\begin{lstlisting}[language=bash, caption=Flujo de comunicación entre servicios]
Usuario (Navegador)
    |
    | HTTP (Puerto 80)
    v
Frontend (Vue 3 + Nginx)
    |
    | REST API (Puerto 3000)
    | Endpoint: http://backend:3000/api
    v
Backend (Node.js + Express)
    |
    | MongoDB Protocol (Puerto 27017)
    | URI: mongodb://mongodb:27017/todoapp
    v
MongoDB (Base de Datos)
\end{lstlisting}

\textbf{Red Docker:} Todos los servicios están conectados a través de una red bridge personalizada llamada \texttt{todo-network}, lo que permite la comunicación interna mediante nombres de servicio.

\textbf{Variables de Entorno:}
\begin{itemize}
    \item Frontend: \texttt{VITE\_API\_URL=http://localhost:3000/api}
    \item Backend: \texttt{MONGODB\_URI=mongodb://mongodb:27017/todoapp}
    \item Backend: \texttt{CORS\_ORIGIN=http://localhost}
\end{itemize}

\subsection{Estructura de carpetas y archivos}

\begin{lstlisting}[caption=Estructura del proyecto]
todo-app-fullstack/
|-- frontend/
|   |-- src/
|   |   |-- components/
|   |   |   |-- TodoForm.vue
|   |   |   |-- TodoItem.vue
|   |   |   |-- TodoFilters.vue
|   |   |   `-- TodoStats.vue
|   |   |-- services/
|   |   |   `-- api.ts
|   |   |-- stores/
|   |   |   `-- todoStore.ts
|   |   |-- types/
|   |   |   `-- index.ts
|   |   |-- App.vue
|   |   `-- main.ts
|   |-- Dockerfile
|   |-- nginx.conf
|   `-- package.json
|
|-- backend/
|   |-- src/
|   |   |-- config/
|   |   |   `-- database.ts
|   |   |-- controllers/
|   |   |   `-- todoController.ts
|   |   |-- middleware/
|   |   |   `-- errorHandler.ts
|   |   |-- models/
|   |   |   `-- Todo.ts
|   |   |-- routes/
|   |   |   `-- todoRoutes.ts
|   |   |-- types/
|   |   |   `-- index.ts
|   |   `-- server.ts
|   |-- Dockerfile
|   |-- .env
|   `-- package.json
|
|-- docs/
|   `-- primer_avance.tex
|
|-- docker-compose.yml
|-- .dockerignore
|-- README.md
`-- SETUP.md
\end{lstlisting}

\newpage

% ===============================
% EVIDENCIA DE EJECUCIÓN
% ===============================
\section{Evidencia de Ejecución}

\subsection{Creación del Repositorio en GitHub}

\textbf{Descripción:} Se creó el repositorio \texttt{todo-app-fullstack} en GitHub con visibilidad pública. El repositorio incluye todo el código fuente, documentación y configuraciones necesarias.

\textbf{URL del repositorio:} \url{https://github.com/LeandroNV/todo-app-fullstack}

\textbf{Estadísticas del primer commit:}
\begin{itemize}
    \item Commit ID: \texttt{6a791e5}
    \item Archivos: 46
    \item Líneas añadidas: 9,172
    \item Mensaje: "feat: Aplicación To-Do Full Stack completa"
\end{itemize}

\textbf{Captura esperada:} Página principal del repositorio en GitHub mostrando los archivos, README y descripción del proyecto.

\subsection{Terminal con docker compose up}

\textbf{Descripción:} Ejecución del comando \texttt{docker compose up} que:
\begin{enumerate}
    \item Construye las imágenes Docker para frontend y backend
    \item Descarga la imagen de MongoDB 7 Alpine
    \item Crea la red \texttt{todo-network}
    \item Inicia los tres contenedores en el orden correcto (MongoDB → Backend → Frontend)
    \item Ejecuta los healthchecks para verificar que los servicios estén funcionando
\end{enumerate}

\textbf{Comandos esperados en la terminal:}
\begin{lstlisting}[language=bash]
$ docker compose up

[+] Building 45.2s (24/24) FINISHED
[+] Running 4/4
 - Network todo-network     Created
 - Container todo-mongodb   Started
 - Container todo-backend   Started
 - Container todo-frontend  Started
\end{lstlisting}

\textbf{Captura esperada:} Terminal mostrando los logs de los tres contenedores ejecutándose simultáneamente, con mensajes de:
\begin{itemize}
    \item MongoDB: "Waiting for connections on port 27017"
    \item Backend: "MongoDB conectado exitosamente" y "Servidor corriendo en http://localhost:3000"
    \item Frontend: Nginx iniciado correctamente
\end{itemize}

\subsection{Navegador accediendo a la aplicación}

\textbf{Descripción:} Aplicación funcionando en \texttt{http://localhost} mostrando:
\begin{itemize}
    \item Dashboard con estadísticas (Total, Completadas, Pendientes, Alta Prioridad, Progreso)
    \item Filtros de búsqueda, estado y prioridad
    \item Formulario para crear nuevas tareas
    \item Indicador de conexión en verde (Conectado)
    \item Interfaz responsive con diseño moderno en Tailwind CSS
\end{itemize}

\textbf{URLs de prueba:}
\begin{itemize}
    \item Frontend: \url{http://localhost}
    \item API Backend: \url{http://localhost:3000/api/todos}
    \item Health Check: \url{http://localhost:3000/health}
\end{itemize}

\textbf{Captura esperada:} Navegador mostrando la interfaz completa de la aplicación con todas las secciones visibles y el indicador "Conectado" en verde.

\subsection{Verificación de contenedores}

\textbf{Comando adicional:}
\begin{lstlisting}[language=bash]
$ docker ps

CONTAINER ID   IMAGE              STATUS                    PORTS
xxx            todo-frontend      Up 2 minutes (healthy)    0.0.0.0:80->80/tcp
xxx            todo-backend       Up 2 minutes (healthy)    0.0.0.0:3000->3000/tcp
xxx            todo-mongodb       Up 2 minutes (healthy)    0.0.0.0:27017->27017/tcp
\end{lstlisting}

\textbf{Captura esperada:} Terminal mostrando los tres contenedores con estado "healthy" y los puertos correctamente mapeados.

\newpage

% ===============================
% PROBLEMAS ENFRENTADOS
% ===============================
\section{Problemas Enfrentados}

Durante el desarrollo de este primer avance, se enfrentaron los siguientes desafíos técnicos:

\subsection{Configuración de CORS}

\textbf{Problema:} Al intentar conectar el frontend con el backend, se presentaban errores de CORS (Cross-Origin Resource Sharing), impidiendo que las peticiones HTTP llegaran al API.

\textbf{Solución:} 
\begin{itemize}
    \item Se configuró el middleware \texttt{cors} en el backend de Express
    \item Se estableció la variable de entorno \texttt{CORS\_ORIGIN} para controlar los orígenes permitidos
    \item En producción con Docker, se ajustó a \texttt{http://localhost} para permitir las peticiones desde el contenedor de frontend
\end{itemize}

\subsection{Comunicación entre contenedores Docker}

\textbf{Problema:} Los contenedores no podían comunicarse entre sí usando \texttt{localhost}, ya que cada contenedor tiene su propia red interna.

\textbf{Solución:}
\begin{itemize}
    \item Se creó una red bridge personalizada (\texttt{todo-network})
    \item Se configuraron los servicios para usar nombres de servicio en lugar de localhost (ejemplo: \texttt{mongodb://mongodb:27017})
    \item Se implementaron healthchecks para asegurar que los servicios estén listos antes de que otros dependan de ellos
\end{itemize}

\subsection{Build multi-stage para el frontend}

\textbf{Problema:} La imagen Docker del frontend era demasiado grande (más de 1GB) al incluir todas las dependencias de desarrollo.

\textbf{Solución:}
\begin{itemize}
    \item Se implementó un Dockerfile multi-stage
    \item Primera etapa: construye la aplicación con Node.js y todas las dependencias
    \item Segunda etapa: copia solo los archivos estáticos compilados a una imagen Nginx Alpine (menos de 50MB)
    \item Resultado: reducción del 95\% en el tamaño de la imagen
\end{itemize}

\subsection{Persistencia de datos en MongoDB}

\textbf{Problema:} Los datos se perdían cada vez que se reiniciaba el contenedor de MongoDB.

\textbf{Solución:}
\begin{itemize}
    \item Se configuraron volúmenes Docker (\texttt{mongodb\_data} y \texttt{mongodb\_config})
    \item Los volúmenes persisten los datos incluso cuando los contenedores se detienen o eliminan
    \item Se documentó el proceso de backup y restauración
\end{itemize}

\subsection{Variables de entorno en tiempo de build}

\textbf{Problema:} Vue 3 con Vite requiere que las variables de entorno estén disponibles en tiempo de compilación, no en tiempo de ejecución.

\textbf{Solución:}
\begin{itemize}
    \item Se configuraron las variables con el prefijo \texttt{VITE\_}
    \item Se pasaron como argumentos en el docker-compose.yml
    \item Se documentó cómo modificarlas para diferentes entornos (desarrollo, staging, producción)
\end{itemize}

\newpage

% ===============================
% PRÓXIMOS PASOS
% ===============================
\section{Próximos Pasos}

Para el siguiente avance del proyecto, se planea implementar las siguientes mejoras y funcionalidades:

\subsection{Integración Continua con Jenkins}

\begin{itemize}
    \item \textbf{Pipeline automatizado:} Configurar un Jenkinsfile con stages para:
    \begin{itemize}
        \item Checkout del código desde GitHub
        \item Instalación de dependencias
        \item Ejecución de linters (ESLint, Prettier)
        \item Compilación de TypeScript
        \item Build de imágenes Docker
        \item Push a Docker Registry
    \end{itemize}
    
    \item \textbf{Webhooks:} Configurar GitHub webhooks para trigger automático del pipeline en cada push o pull request
    
    \item \textbf{Notificaciones:} Implementar notificaciones de éxito o fallo del pipeline (email, Slack, etc.)
\end{itemize}

\subsection{Pruebas Automatizadas}

\begin{itemize}
    \item \textbf{Tests unitarios:}
    \begin{itemize}
        \item Frontend: Vitest + Vue Test Utils
        \item Backend: Jest
        \item Objetivo: mínimo 80\% de cobertura de código
    \end{itemize}
    
    \item \textbf{Tests de integración:}
    \begin{itemize}
        \item Supertest para endpoints del API
        \item Validar respuestas y códigos de estado HTTP
        \item Probar flujos completos CRUD
    \end{itemize}
    
    \item \textbf{Tests End-to-End (E2E):}
    \begin{itemize}
        \item Playwright o Cypress
        \item Simular interacciones de usuario reales
        \item Validar flujos completos de la aplicación
    \end{itemize}
\end{itemize}

\subsection{Despliegue Continuo (CD)}

\begin{itemize}
    \item \textbf{Ambientes:} Configurar tres entornos separados:
    \begin{itemize}
        \item Development (rama develop)
        \item Staging (rama staging)
        \item Production (rama main)
    \end{itemize}
    
    \item \textbf{Estrategias de despliegue:}
    \begin{itemize}
        \item Blue-Green deployment
        \item Rolling updates
        \item Rollback automático en caso de fallo
    \end{itemize}
    
    \item \textbf{Docker Registry:} Configurar un registry privado para almacenar las imágenes
\end{itemize}

\subsection{Monitoreo y Observabilidad}

\begin{itemize}
    \item \textbf{Logging:}
    \begin{itemize}
        \item Centralizar logs con ELK Stack (Elasticsearch, Logstash, Kibana)
        \item Logs estructurados en formato JSON
        \item Niveles de log apropiados (error, warn, info, debug)
    \end{itemize}
    
    \item \textbf{Métricas:}
    \begin{itemize}
        \item Prometheus para recolección de métricas
        \item Grafana para visualización
        \item Alertas configuradas para errores críticos
    \end{itemize}
    
    \item \textbf{Healthchecks:} Endpoints de health mejorados con información detallada del estado de cada servicio
\end{itemize}

\subsection{Mejoras de Seguridad}

\begin{itemize}
    \item \textbf{Autenticación:} Implementar JWT para autenticación de usuarios
    \item \textbf{Autorización:} Sistema de roles y permisos
    \item \textbf{HTTPS:} Configurar certificados SSL/TLS
    \item \textbf{Secrets:} Usar Docker secrets o vault para gestionar credenciales
    \item \textbf{Escaneo de vulnerabilidades:} Integrar Trivy o Snyk en el pipeline
\end{itemize}

\subsection{Optimizaciones}

\begin{itemize}
    \item \textbf{Caché:} Implementar Redis para caché de datos frecuentes
    \item \textbf{CDN:} Servir assets estáticos desde un CDN
    \item \textbf{Database indexing:} Optimizar índices de MongoDB basados en queries más frecuentes
    \item \textbf{Lazy loading:} Implementar carga diferida en componentes Vue
\end{itemize}

\newpage

% ===============================
% CONCLUSIONES
% ===============================
\section{Conclusiones}

En este primer avance se logró:

\begin{itemize}
    \item Desarrollar una aplicación full-stack funcional y completa
    \item Implementar una arquitectura de microservicios con Docker
    \item Establecer las bases para CI/CD con documentación exhaustiva
    \item Crear un repositorio bien estructurado y documentado en GitHub
    \item Configurar comunicación eficiente entre los servicios
\end{itemize}

El proyecto está ahora en un estado sólido para continuar con la implementación de CI/CD, pruebas automatizadas y mejoras de seguridad en los siguientes avances.

\vspace{1cm}

\textbf{Stack tecnológico final:}
\begin{itemize}
    \item Frontend: Vue 3, TypeScript, Tailwind CSS, Vite, Nginx
    \item Backend: Node.js 20, Express, TypeScript, Mongoose
    \item Database: MongoDB 7 Alpine
    \item Infraestructura: Docker, Docker Compose
    \item Próximo: Jenkins, Tests automatizados
\end{itemize}

\end{document}

